\section{Results and Conclusions}
\label{sec:conclusions}

\textbf{Main question.} Combined, municipal socioeconomics and school-level
resources predict Bagrut outcomes usefully: nested out-of-fold $R^2$ of 0.40 to
0.56 (Table~\ref{tab:leaderboard}) with near-zero mean residuals
(Figure~\ref{fig:residuals}). \emph{Combined} carries it: municipal data alone
reaches only 0.06 to 0.23, and the school-level features lift every target by
$+0.309$ on average, in every outer fold (Table~\ref{tab:ablation}).

\textbf{Secondary question A.} Mathematics is the more resilient subject
(Table~\ref{tab:secondary}): the cluster~2 to cluster~9 gap is smaller for Math
in achievement and about a third of English's in participation. Overachieving
schools clearly exist, 87 of 460 low-cluster schools, with significantly higher
advanced-track participation than their peers.

\textbf{Secondary question B.} The strongest school-level predictor is the
Ministry's per-school nurture index, first for all four targets, followed by the
transport-related budget field, school size, and class size. Five per-student
budget ratios are confirmed for every target (Section~\ref{sec:importance}).
Municipal cluster never exceeds 5.3\% of attribution and is dropped for English
average grade.

\textbf{Conclusion.} Socioeconomic status matters, but the municipal average is
much coarser than the school-level measure. Measured per school the same
construct dominates every model. The town-level cluster is among the weakest
features retained and is rejected entirely for one target.

\textbf{Limitations and further work.} The institutional data is a single 2014/15
snapshot applied to all four exam years, so within-school change cannot be
separated. Suppressed grade cells are small cohorts, so the smallest schools are
underrepresented, and complete-case modelling removes a further 9--12.5\% of
rows. The transport field's negative values leave its meaning unconfirmed.
Results are associational, and this model is not suitable for individual
admissions decisions or for withholding resources from schools. Natural
extensions: per-year budget records to test whether spending \emph{precedes}
outcomes, a study of the 87 overachievers, and bounded-response models for
participation rates.
